% !TEX program = xelatex
% 电像法与泊松核
\documentclass[fontset=windows, chinese, 10pt, xcolor=dvipsnames]{beamer}
% Shared preamble for the chapter 10 standalone lectures.
% Each lecture sets \title and \subtitle after this file is read.

\usepackage{amsmath,amssymb,mathtools}
\usepackage{xeCJK}
\usepackage{zhnumber}
\usepackage{tikz}
\usepackage{tikz-3dplot}
\usetikzlibrary{arrows.meta,calc,positioning}

\usetheme{metropolis}
\usepackage{beamercolorthemeHistNanjing}

\newcommand{\red}[1]{\textcolor{red}{#1}}
\newcommand{\blue}[1]{\textcolor{blue}{#1}}
\newcommand{\green}[1]{\textcolor{green}{#1}}
\newcommand{\gray}[1]{\textcolor{gray}{#1}}
\newcommand{\orange}[1]{\textcolor{orange}{#1}}

\newcommand{\sssub}[1]{ {\mathrm{#1} }}  % script size subscript\\\
\newcommand{\funcd}[2]{ \frac{d{#1}}{d{#2}}} %function full derivative
\newcommand{\funcpd}[2]{ \frac{\partial{#1}}{\partial{#2}}} %function partial derivative
\newcommand{\fnald}[2] { \frac{\delta {#1}}{\delta {#2}}} % functional derivative

\renewcommand{\Re}{\operatorname{Re}}
\renewcommand{\Im}{\operatorname{Im}}
\def\ext{\sssub{ext}}
\def\x{\sssub{x}}
\def\xc{\sssub{xc}}
\def\H{\sssub{H}}
\def\Hx{\sssub{Hx}}
\def\Hxc{\sssub{Hxc}}
\def\c{\sssub{c}}
\def\s{\sssub{s}}

% \def\ext{\sssub{ext}}
% \def\ext{\sssub{ext}}
% \def\ext{\sssub{ext}}
% \def\ext{\sssub{ext}}
% \defsssub{ext}


% definitions of bold faced symbols
\def\brpppp{{\mathbf{r}^{\prime\prime\prime\prime}}}
\def\brppp{{\mathbf{r}^{\prime\prime\prime}}}
\def\brpp{{\mathbf{r}^{\prime\prime}}}
\def\brp{{\mathbf{r}^{\prime}}}
\def\bzp{{\mathbf{z}^{\prime}}}
\def\bxp{{\mathbf{x}^{\prime}}}
\def\tp{{{t}^{\prime}}}
\def\tpp{{{t}^{\prime\prime}}}
\def\tppp{{{t}^{\prime\prime\prime}}}

\def\tbr{{\tilde{\mathbf{r}}}}
\def\bk{{\mathbf{k}}}
\def\br{{\mathbf{r}}}
\def\bp{{\mathbf{p}}}
\def\bq{{\mathbf{q}}}
\def\bv{{\mathbf{v}}}
\def\bz{{\mathbf{z}}}
\def\mbf{{\mathbf{f}}}

\def\bx{{\mathbf{x}}}
\def\bR{{\mathbf{R}}}
\def\bM{{\mathbf{M}}}
\def\bP{{\mathbf{P}}}
\def\bT{{\mathbf{T}}}
\def\bK{{\mathbf{K}}}
\def\bA{{\mathbf{A}}}
\def\bB{{\mathbf{B}}}
\def\bD{{\mathbf{D}}}
\def\bE{{\mathbf{E}}}
\def\bF{{\mathbf{F}}}
\def\bG{{\mathbf{G}}}
\def\bH{{\mathbf{H}}}
\def\bI{{\mathbf{I}}}
\def\bJ{{\mathbf{J}}}
\def\bP{{\mathbf{P}}}
\def\bS{{\mathbf{S}}}
\def\bV{{\mathbf{V}}}
\def\bX{{\mathbf{X}}}
\def\bY{{\mathbf{Y}}}
\def\dd{{\mathop{}\!\mathrm{d}}}
\def\dr{{\!\!\dd\br\,}}
\def\dx{{\!\!\dd\bx\,}}
\def\dy{{\!\!\dd\by\,}}
\def\dk{{\!\!\dd\bk\,}}
\def\dG{{\!\!\dd\bG\,}}


\def\half{{\frac{1}{2}}}
\def\tr{{\mathrm{tr}}}


\def\Arg{{\mathrm{Arg}\;}}
\def\arg{{\mathrm{arg}\;}}
\def\Res{{\mathrm{Res}\;}}

% unit-vector basis (bold under XeLaTeX; \mathbf can look non-bold)
\newcommand{\ue}{\boldsymbol{e}}
\newcommand{\un}{\hat{\boldsymbol{n}}}


\author{罗凯}
\institute{南京理工大学物理学院}
\date{\zhdigits{\the\year},\zhganzhinian{\the\year}年秋}

\usebackgroundtemplate{%
\begin{tikzpicture}[remember picture,overlay]
  \node[opacity=0.05, at=(current page.center)] {
    \includegraphics[width=0.6\paperwidth]{njust-logo.jpeg}
  };
\end{tikzpicture}}


\title{格林函数法}
\subtitle{第二讲：电像法}

\begin{document}

\begin{frame}
  \titlepage
\end{frame}

\begin{frame}{本讲提纲}
  边界简单时，格林函数写成自由空间的基本解加上区域外面的像。
  \begin{enumerate}
    \item 三维、二维的基本解；
    \item 接地球和圆盘的反演像；
    \item 四个泊松核：球、半空间、圆盘、半平面；
    \item 圆盘上 $A\cos\phi$ 回到柱坐标的分离变量解。
  \end{enumerate}
\end{frame}

\section{基本解}

\begin{frame}{奇异部分与调和修正}
  \[
    G=G_0+G_1,
    \qquad
    \nabla^2 G_0=\delta(\mathbf{r}-\mathbf{r}_0),
    \qquad
    \nabla^2 G_1=0.
  \]
  狄利克雷问题还要求 $G_1|_{\Sigma}=-G_0|_{\Sigma}$。
  \vspace{0.4em}
  像源放在区域外面。区域内部仍然只有原来的点源，边界条件靠像去配。
\end{frame}

\begin{frame}{三维}
  源在原点时，径向解是 $-C_1/r+C_2$。无穷远处趋于零，则 $C_2=0$。对小球用散度定理，
  \[
    \iint\frac{\partial G_0}{\partial r}\,\mathrm{d}S=4\pi C_1=1,
  \]
  所以 $C_1=1/(4\pi)$。源点移到 $\mathbf{r}_0$，
  \begin{equation*}
    \boxed{\displaystyle
      G_0=-\frac{1}{4\pi|\mathbf{r}-\mathbf{r}_0|}.}
  \end{equation*}
\end{frame}

\begin{frame}{二维}
  同样的手续给出
  \begin{equation*}
    \boxed{\displaystyle
      G_0=\frac{1}{2\pi}\ln|\mathbf{r}-\mathbf{r}_0|.}
  \end{equation*}
  静电里，方程是 $\nabla^2 V=-\rho/\varepsilon_0$ 时，这两个基本解对应自由空间中电荷 $-\varepsilon_0$ 的电位，电荷要改号。
\end{frame}

\section{像}

\begin{frame}{接地球}
  球半径 $a$，源在 $r_0<a$。像放在同一射线上的反演点，
  \[
    \mathbf{r}_1=\frac{a^2}{r_0^2}\mathbf{r}_0,
    \qquad r_1=\frac{a^2}{r_0}>a.
  \]
  球面上 $|\mathbf{r}|=a$ 时，$|\mathbf{r}-\mathbf{r}_1|=(a/r_0)\,|\mathbf{r}-\mathbf{r}_0|$。像的强度取 $a/r_0$，两项相消，
  \begin{equation*}
    G
    =-\frac{1}{4\pi|\mathbf{r}-\mathbf{r}_0|}
     +\frac{a/r_0}{4\pi\bigl|\mathbf{r}-(a^2/r_0^2)\mathbf{r}_0\bigr|}.
  \end{equation*}
  像在球外，球内 $\nabla^2 G$ 只有原来的 $\delta$。
\end{frame}

\begin{frame}{圆盘}
  $\rho<a$ 时用同一个反演点，
  \[
    G
    =\frac{1}{2\pi}\ln|\mathbf{r}-\mathbf{r}_0|
     -\frac{1}{2\pi}\ln|\mathbf{r}-\mathbf{r}_1|
     +\frac{1}{2\pi}\ln\frac{a}{r_0}.
  \]
  常数不破坏调和性。圆周上两个距离相差因子 $a/r_0$，对数相加是零。
\end{frame}

\section{泊松核}

\begin{frame}{球的泊松积分}
  $f=0$ 时只剩 $\partial G/\partial n_0$。在 $r_0=a$ 上像点和源点重合，但法向导数不是零，
  \[
    \frac{\partial G}{\partial r_0}\Big|_{r_0=a}
    =\frac{1}{4\pi a}\frac{a^2-r^2}{R^3},
  \]
  $R^2=r^2+a^2-2ar\cos\Theta$。乘上面积元 $a^2\sin\theta_0\,\mathrm{d}\theta_0\,\mathrm{d}\phi_0$，
  \begin{equation*}
    u
    =\frac{a}{4\pi}\int
      f(\theta_0,\phi_0)
      \frac{a^2-r^2}{(a^2-2ar\cos\Theta+r^2)^{3/2}}
      \,\mathrm{d}\Omega_0.
  \end{equation*}
\end{frame}

\begin{frame}{球心是平均值}
  $r\to a$ 时核变成球面上的 $\delta$，$u$ 回到边界值。
  \vspace{0.4em}
  球心处 $\cos\Theta$ 消失，分母是 $a^3$，分子是 $a^2$，
  \[
    u(0)=\frac{1}{4\pi}\int f\,\mathrm{d}\Omega.
  \]
  调和函数在球心的值等于边界值在球面上的平均。
\end{frame}

\begin{frame}{上半空间}
  像与源反号，放在 $(x_0,y_0,-z_0)$。平面 $z=0$ 上两个距离相等，$G=0$。
  \vspace{0.3em}
  上半空间的外法向是 $-\mathbf{e}_z$。平面上
  \[
    \frac{\partial G}{\partial n_0}
    =\frac{z}{2\pi\bigl[(x-x_0)^2+(y-y_0)^2+z^2\bigr]^{3/2}}.
  \]
  因此
  \[
    u(x,y,z)
    =\frac{z}{2\pi}\iint
      \frac{f(x_0,y_0)}
           {\bigl[(x-x_0)^2+(y-y_0)^2+z^2\bigr]^{3/2}}
      \,\mathrm{d}x_0\,\mathrm{d}y_0.
  \]
\end{frame}

\begin{frame}{圆盘与半平面}
  圆盘上对外法向 $\partial/\partial\rho_0$ 求导，在 $\rho_0=a$ 处
  \[
    u(\rho,\phi)
    =\frac{a^2-\rho^2}{2\pi}\int_0^{2\pi}
      \frac{f(\phi_0)}
           {a^2-2a\rho\cos(\phi-\phi_0)+\rho^2}
      \,\mathrm{d}\phi_0.
  \]
  上半平面用反号像点 $(x_0,-y_0)$，外法向 $-\mathbf{e}_y$，
  \[
    u(x,y)
    =\frac{y}{\pi}\int_{-\infty}^{\infty}
      \frac{f(x_0)}{(x-x_0)^2+y^2}\,\mathrm{d}x_0.
  \]
  $f=1$ 时两个积分都等于 $1$。常值函数调和，边界值也是 $1$。
\end{frame}

\begin{frame}{圆盘上的 $A\cos\phi$}
  边界 $u(a,\phi)=A\cos\phi$。内部有限的解是 $\rho^{n}$ 的傅里叶级数，边界只有 $\cos\phi$，所以
  \[
    u(\rho,\phi)=A\frac{\rho}{a}\cos\phi.
  \]
  这就是柱坐标一章圆内边值的那个函数。泊松积分把分离变量的结果包进一个核对边界值的积分里。
\end{frame}

\begin{frame}{这一讲留下的公式}
  \[
    G_0^{(3)}=-\frac{1}{4\pi R},
    \qquad
    G_0^{(2)}=\frac{1}{2\pi}\ln R.
  \]
  接地球的像在反演点 $r_1=a^2/r_0$，强度是 $a/r_0$。
  \vspace{0.35em}
  球的泊松核是 $(a^2-r^2)/R^3$，半空间是 $z/R^3$，圆盘是 $(a^2-\rho^2)/R^2$，半平面是 $y/R^2$。
\end{frame}

\end{document}
